% ============ 第5章 §5.6 Schwartz 分布的卷积 (原书 pp.120-122) ============

\section{Schwartz 分布的卷积}\label{sec:5.6}
\index{0symbolSconv@$S(\mathbb{R}^{n})$ 与 $S'(\mathbb{R}^{n})$ 的卷积}




作为属于 $\mathcal{S}(\mathbb{R}^{n})$ 的函数之间卷积的定义的动机, 如下定义 $f \in \mathcal{S}(\mathbb{R}^{n})$ 与 $g \in \mathcal{S}'(\mathbb{R}^{n})$ 的卷积.

\begin{definition}[$\mathcal{S}(\mathbb{R}^{n})$ 与 $\mathcal{S}'(\mathbb{R}^{n})$ 的卷积]\label{def:5.7}
对 $f \in \mathcal{S}(\mathbb{R}^{n})$ 与 $g \in \mathcal{S}'(\mathbb{R}^{n})$, 函数 $f * g$ 定义为
\[
f * g(x) \equiv \langle g, f(x - \cdot)\rangle \quad (x \in \mathbb{R}^{n}).
\]
\end{definition}

在定义~\ref{def:5.7} 中, 即使 $g \in \mathcal{S}'(\mathbb{R}^{n})$ 也能取到值 $f * g(x)$, 这是重要的点. 对 $f \in L^{1}(\mathbb{R}^{n})$ 的 $F_{f}$ 在例~\ref{exa:5.9} 中给出. 汇总一下这个运算的基本性质.

\begin{theorem}[卷积的基本性质]\label{thm:5.12}
设 $f, g \in \mathcal{S}(\mathbb{R}^{n})$, $h \in \mathcal{S}'(\mathbb{R}^{n})$.
(1) 对 $x \in \mathbb{R}^{n}$, $L^{1}$ 函数 f 与 $L^{1}$ 函数 g 的卷积 $f * g(x) = \int_{\mathbb{R}^{n}}f(x-y)g(y)\dd y$ 与定义~\ref{def:5.7} 的 $F_{f} * g(x)$ 一致.
(2) $f * h \in C^{\infty}(\mathbb{R}^{n})$. 并且, 对 $\alpha \in \mathbb{N}_{0}^{n}$,
\begin{equation}\tag{5.18}
\partial^{\alpha}[f*h] = (\partial^{\alpha}f)*h = f*\partial^{\alpha}h
\label{eq:5.18}
\end{equation}
成立. 又, 对所有 $\alpha \in \mathbb{N}_{0}^{n}$, 存在 $N = N_{\alpha} \in \mathbb{N}_{0}$, 使得
\begin{equation}\tag{5.19}
\sup_{x \in \mathbb{R}^{n}}|\partial^{\alpha}[f*h](x)| \lesssim (1+|x|)^{N}
\label{eq:5.19}
\end{equation}
成立. 因此, 特别地, $f*h = F_{f*h} \in \mathcal{S}'(\mathbb{R}^{n})$ 可以这样看待.
(3) 结合律 $(f*g)*h = f*(g*h)$ 成立.
\end{theorem}

\begin{proof}
(1) 由 $f \in L^{1}(\mathbb{R}^{n}) \mapsto F_{f} \in \mathcal{S}'(\mathbb{R}^{n})$ 与卷积的定义显然.

(2) 只要证明
\begin{equation}\tag{5.20}
\partial^{\alpha}[f*h] = f*\partial^{\alpha}h, |\alpha| = 1
\label{eq:5.20}
\end{equation}
就可以反复使用定义~\ref{def:5.6} 与式~\eqref{eq:5.20} 来证明, 所以证明式~\eqref{eq:5.20} 就足够了. $f*h \in C^{\infty}(\mathbb{R}^{n})$ 的证明使用例~\ref{exa:5.10}. 若 $|\alpha| = 1$, 即存在 $j = 1, 2, \ldots$ 使得 $\partial^{\alpha} = \partial_{x_{j}}$, 则式~\eqref{eq:5.18} 由分布的微分表示为差分这一点是显然的. 如上所述, $|\alpha| \geq 2$ 的情形可以用数学归纳法证明式~\eqref{eq:5.18}.

为了证明不等式~\eqref{eq:5.19}, 由定理~\ref{thm:5.3} 首先注意存在自然数 N, 使得对所有 $\varphi \in \mathcal{S}(\mathbb{R}^{n})$,
\begin{equation}\tag{5.21}
|\langle h, \varphi\rangle| \leq Np_{N}(\varphi)
\label{eq:5.21}
\end{equation}
成立. 在式~\eqref{eq:5.21} 中取 $f(x - \cdot)$ 作为 $\varphi$. 由 $p_{N}$ 的定义, 通过简单计算可知 $p_{N}(\partial^{\alpha}f(x - \cdot)) \lesssim (1+|x|)^{N}$, 所以 $|\partial^{\alpha}[f*h](x)| \lesssim (1+|x|)^{N}$ 随之成立.

(3) 由于
\[
f*g = \sum_{k \in \mathbb{Z}^{n}}\int_{k+[0,1]^{n}}g(\cdot - y)f(y)\dd y
\]
在 $\mathcal{S}(\mathbb{R}^{n})$ 的拓扑中成立, 固定 $x \in \mathbb{R}^{n}$ 后,
\[
(f*g)*h(x) = \sum_{k \in \mathbb{Z}^{n}}\left\langle h, \int_{k+[0,1]^{n}}g(x - \cdot - y)f(y)\dd y\right\rangle.
\]
由引理~\ref{lem:5.2}, 将积分改写为 Riemann 和风格的操作在 $\mathcal{S}(\mathbb{R}^{n})$ 的拓扑中也有效, 所以
\[
(f*g)*h(x) = \sum_{k \in \mathbb{Z}^{n}}\left(\int_{k+[0,1]^{n}}f(y)g*h(x-y)\dd y\right).
\]
由于被积函数 $f(\cdot)g*h(x - \cdot)$ 在 $\mathbb{R}^{n}$ 上可积,
\[
(f*g)*h(x) = \int_{\mathbb{R}^{n}}f(y)g*h(x-y)\dd y.
\]
由以上, 证明了 $f*(g*h) = (f*g)*h$.
\end{proof}

\begin{exercise}\label{exer:5.12}
当 $f \in \mathcal{S}(\mathbb{R}^{n})$, $g \in \mathcal{S}'(\mathbb{R}^{n})$ 时, 证明 $\mathcal{F}[f*g] = (2\pi)^{\frac{n}{2}}\mathcal{F}g \cdot \mathcal{F}f$. 但是, 对 $F \in \mathcal{S}'(\mathbb{R}^{n})$, $\Phi \in \mathcal{S}(\mathbb{R}^{n})$, 通过 $\langle F \cdot \Phi, \varphi\rangle = \langle F, \Phi\varphi\rangle$ 定义 $F \cdot \Phi \in \mathcal{S}'(\mathbb{R}^{n})$.
\end{exercise}
